% --- Archon physics formalization source begin ---
% archon:physics
% archon:covers PhyXMiniProblems/problem_phyx_mini_0025.lean
% archon:source-report reports/phyx_mini/problem_phyx_mini_0025.source.json

\chapter{Physics problem phyx\_mini\_0025}
\label{ch:PhyXMiniProblems_problem_phyx_mini_0025}

\paragraph{Problem source.}
\#\# Physical scenario

Figure shows an overhead view of a room of square floor area and side \textbackslash{}( L \textbackslash{}). At the center of the room is a mirror set in a vertical plane and rotating on a vertical shaft at angular speed \textbackslash{}( \textbackslash{}omega \textbackslash{}) about an axis coming out of the page. A bright red laser beam enters from the center point on one wall of the room and strikes the mirror. As the mirror rotates, the reflected laser beam creates a red spot sweeping across the walls of the room.

\#\# Question

In what time interval does the spot change from its minimum to its maximum speed?

\#\# Answer choices

- A: A: \textbackslash{}( \textbackslash{}frac\{\textbackslash{}pi\}\{4\textbackslash{}omega\} \textbackslash{})

- B: B: \textbackslash{}( \textbackslash{}frac\{\textbackslash{}pi\}\{12\textbackslash{}omega\} \textbackslash{})

- C: C: \textbackslash{}( \textbackslash{}frac\{\textbackslash{}pi\}\{6\textbackslash{}omega\} \textbackslash{})

- D: D: \textbackslash{}( \textbackslash{}frac\{\textbackslash{}pi\}\{8\textbackslash{}omega\} \textbackslash{})

\#\# Figure caption (auxiliary; use the image as primary evidence)

The image depicts a diagram involving a mirror within a square enclosure with side length labeled \textbackslash{}( L \textbackslash{}).

- At the center of the diagram, there is a mirror, tilted at an angle, shown in light blue. It is represented to be rotating counterclockwise, indicated by an arrow and the angular velocity symbol \textbackslash{}( \textbackslash{}omega \textbackslash{}).

- Two red arrows represent beams or rays reflecting off the mirror: one incoming vertically from the bottom and another reflecting at an angle upwards towards the right. 

- A dashed horizontal line passes through point \textbackslash{}( O \textbackslash{}) slightly to the right of the mirror, with a vertical measurement labeled \textbackslash{}( x \textbackslash{}) extending from this line to the point \textbackslash{}( O \textbackslash{}).

- The letters \textbackslash{}( L \textbackslash{}) on the top and left sides of the square indicate the side length of the enclosure.

- The structure of the diagram suggests it involves optical principles related to reflection and possibly the rotation of the mirror affecting the beam's path.

\#\# Dataset metadata

Geometrical Optics; Physical Model Grounding Reasoning, Predictive Reasoning

\paragraph{Recorded answer/context.}
D: D: \textbackslash{}( \textbackslash{}frac\{\textbackslash{}pi\}\{8\textbackslash{}omega\} \textbackslash{})

\paragraph{Figure/image path.}
/root/proposal\_for\_physic/science-mango/phyx\_mini\_run/phyx\_data/test\_image/25.png

\paragraph{Formalization target.}
create a compiling Lean file with sorry bodies at `PhyXMiniProblems/problem\_phyx\_mini\_0025.lean`. The Lean declarations must preserve the physical quantities, dimensions or dimensional roles, figure labels, governing-law hypotheses, and final relation expressed by this problem.
Use Mathlib/Physlib names found through LeanExplore where available. If a physics API is missing, introduce faithful local abstractions rather than scalar placeholder aliases.

\begin{theorem}[Physics formalization target]
\label{thm:physics:phyx_mini_0025:target}
\lean{PhyXMiniProblems.ProblemPhyXMini0025.spotMinimumToMaximumSpeedInterval}
\uses{def:physics:phyx-mini-0025:phyxminiproblems-problemphyxmini0025-dimtime, def:physics:phyx-mini-0025:phyxminiproblems-problemphyxmini0025-rotatingmirrorexperiment, def:physics:phyx-mini-0025:phyxminiproblems-problemphyxmini0025-matchesrotatingmirrorfigure, def:physics:phyx-mini-0025:phyxminiproblems-problemphyxmini0025-hasphysicalparameters, def:physics:phyx-mini-0025:phyxminiproblems-problemphyxmini0025-rotatingmirroropticslaws, def:physics:phyx-mini-0025:phyxminiproblems-problemphyxmini0025-iseastwallsweepfromotocorner, def:physics:phyx-mini-0025:phyxminiproblems-problemphyxmini0025-isminimumspotspeedon, def:physics:phyx-mini-0025:phyxminiproblems-problemphyxmini0025-ismaximumspotspeedon}
The assigned autoformalize agent should translate this physics problem into Lean declarations in the covered file, with theorem and lemma proof bodies written as `by sorry`.
\end{theorem}
\begin{proof}
This is an autoformalization task, not a proof task. Produce statements that can later be proved by the physics prover without weakening the physical model.
\end{proof}

% --- Archon declaration topology begin ---
\section{Declaration topology}
\begin{definition}[Floor Plane]
  \label{def:physics:phyx-mini-0025:phyxminiproblems-problemphyxmini0025-floorplane}
  \lean{PhyXMiniProblems.ProblemPhyXMini0025.FloorPlane}
  The horizontal plane of the overhead room diagram
\end{definition}
\begin{proof}
  This is the declared model definition of Floor Plane.
\end{proof}
\begin{definition}[Dim Length]
  \label{def:physics:phyx-mini-0025:phyxminiproblems-problemphyxmini0025-dimlength}
  \lean{PhyXMiniProblems.ProblemPhyXMini0025.DimLength}
  A physical length, independent of a choice of units
\end{definition}
\begin{proof}
  This is the declared model definition of Dim Length.
\end{proof}
\begin{definition}[Dim Time]
  \label{def:physics:phyx-mini-0025:phyxminiproblems-problemphyxmini0025-dimtime}
  \lean{PhyXMiniProblems.ProblemPhyXMini0025.DimTime}
  A physical time coordinate or time interval
\end{definition}
\begin{proof}
  This is the declared model definition of Dim Time.
\end{proof}
\begin{definition}[Dim Angular Speed]
  \label{def:physics:phyx-mini-0025:phyxminiproblems-problemphyxmini0025-dimangularspeed}
  \lean{PhyXMiniProblems.ProblemPhyXMini0025.DimAngularSpeed}
  A positive-or-negative angular-speed readout, with radians dimensionless
\end{definition}
\begin{proof}
  This is the declared model definition of Dim Angular Speed.
\end{proof}
\begin{definition}[Dim Linear Speed]
  \label{def:physics:phyx-mini-0025:phyxminiproblems-problemphyxmini0025-dimlinearspeed}
  \lean{PhyXMiniProblems.ProblemPhyXMini0025.DimLinearSpeed}
  A physical linear speed along a wall
\end{definition}
\begin{proof}
  This is the declared model definition of Dim Linear Speed.
\end{proof}
\begin{definition}[Floor Point]
  \label{def:physics:phyx-mini-0025:phyxminiproblems-problemphyxmini0025-floorpoint}
  \lean{PhyXMiniProblems.ProblemPhyXMini0025.FloorPoint}
  \uses{def:physics:phyx-mini-0025:phyxminiproblems-problemphyxmini0025-floorplane}
  A point in the floor plane whose coordinates carry length dimension
\end{definition}
\begin{proof}
  This is the declared model definition of Floor Point.
\end{proof}
\begin{definition}[Beam Direction]
  \label{def:physics:phyx-mini-0025:phyxminiproblems-problemphyxmini0025-beamdirection}
  \lean{PhyXMiniProblems.ProblemPhyXMini0025.BeamDirection}
  \uses{def:physics:phyx-mini-0025:phyxminiproblems-problemphyxmini0025-floorplane}
  A direction of propagation, modulo multiplication by a positive scalar
\end{definition}
\begin{proof}
  This is the declared model definition of Beam Direction.
\end{proof}
\begin{definition}[Wall Label]
  \label{def:physics:phyx-mini-0025:phyxminiproblems-problemphyxmini0025-walllabel}
  \lean{PhyXMiniProblems.ProblemPhyXMini0025.WallLabel}
  The four walls of the square, named in the orientation of the primary image
\end{definition}
\begin{proof}
  This is the declared model definition of Wall Label.
\end{proof}
\begin{definition}[Directed Beam]
  \label{def:physics:phyx-mini-0025:phyxminiproblems-problemphyxmini0025-directedbeam}
  \lean{PhyXMiniProblems.ProblemPhyXMini0025.DirectedBeam}
  \uses{def:physics:phyx-mini-0025:phyxminiproblems-problemphyxmini0025-floorpoint, def:physics:phyx-mini-0025:phyxminiproblems-problemphyxmini0025-beamdirection}
  A directed geometrical-optics ray with a dimensioned point of origin
\end{definition}
\begin{proof}
  This is the declared model definition of Directed Beam.
\end{proof}
\begin{definition}[Point On Beam]
  \label{def:physics:phyx-mini-0025:phyxminiproblems-problemphyxmini0025-pointonbeam}
  \lean{PhyXMiniProblems.ProblemPhyXMini0025.PointOnBeam}
  \uses{def:physics:phyx-mini-0025:phyxminiproblems-problemphyxmini0025-floorpoint, def:physics:phyx-mini-0025:phyxminiproblems-problemphyxmini0025-directedbeam}
  A dimensioned point is on the forward half-ray of a directed beam
\end{definition}
\begin{proof}
  This is the declared model definition of Point On Beam.
\end{proof}
\begin{definition}[Square Room Figure]
  \label{def:physics:phyx-mini-0025:phyxminiproblems-problemphyxmini0025-squareroomfigure}
  \lean{PhyXMiniProblems.ProblemPhyXMini0025.SquareRoomFigure}
  \uses{def:physics:phyx-mini-0025:phyxminiproblems-problemphyxmini0025-dimlength, def:physics:phyx-mini-0025:phyxminiproblems-problemphyxmini0025-floorpoint, def:physics:phyx-mini-0025:phyxminiproblems-problemphyxmini0025-walllabel}
  The square-room data and the labels used by the image. The predicate onWall and the east-wall offset are physical incidence/readout interfaces; their metric properties are stated separately in MatchesSquareRoomFigure
\end{definition}
\begin{proof}
  This is the declared model definition of Square Room Figure.
\end{proof}
\begin{definition}[Matches Square Room Figure]
  \label{def:physics:phyx-mini-0025:phyxminiproblems-problemphyxmini0025-matchessquareroomfigure}
  \lean{PhyXMiniProblems.ProblemPhyXMini0025.MatchesSquareRoomFigure}
  \uses{def:physics:phyx-mini-0025:phyxminiproblems-problemphyxmini0025-squareroomfigure}
  Metric and incidence facts read from the square figure. In particular, O is the midpoint of the east wall and its separation from both the center and the northeast corner is L / 2
\end{definition}
\begin{proof}
  This is the declared model definition of Matches Square Room Figure.
\end{proof}
\begin{definition}[Rotating Mirror Experiment]
  \label{def:physics:phyx-mini-0025:phyxminiproblems-problemphyxmini0025-rotatingmirrorexperiment}
  \lean{PhyXMiniProblems.ProblemPhyXMini0025.RotatingMirrorExperiment}
  \uses{def:physics:phyx-mini-0025:phyxminiproblems-problemphyxmini0025-dimtime, def:physics:phyx-mini-0025:phyxminiproblems-problemphyxmini0025-dimangularspeed, def:physics:phyx-mini-0025:phyxminiproblems-problemphyxmini0025-dimlinearspeed, def:physics:phyx-mini-0025:phyxminiproblems-problemphyxmini0025-floorpoint, def:physics:phyx-mini-0025:phyxminiproblems-problemphyxmini0025-directedbeam, def:physics:phyx-mini-0025:phyxminiproblems-problemphyxmini0025-squareroomfigure}
  Time-dependent quantities in the experiment. reflectedAngleFromEastNormal is an unwrapped radian angle: it is zero along the ray from the center to O and increases toward the northeast corner. spotSpeed is the nonnegative linear speed of the spot along the wall currently struck
\end{definition}
\begin{proof}
  This is the declared model definition of Rotating Mirror Experiment.
\end{proof}
\begin{definition}[Matches Rotating Mirror Figure]
  \label{def:physics:phyx-mini-0025:phyxminiproblems-problemphyxmini0025-matchesrotatingmirrorfigure}
  \lean{PhyXMiniProblems.ProblemPhyXMini0025.MatchesRotatingMirrorFigure}
  \uses{def:physics:phyx-mini-0025:phyxminiproblems-problemphyxmini0025-dimtime, def:physics:phyx-mini-0025:phyxminiproblems-problemphyxmini0025-pointonbeam, def:physics:phyx-mini-0025:phyxminiproblems-problemphyxmini0025-matchessquareroomfigure, def:physics:phyx-mini-0025:phyxminiproblems-problemphyxmini0025-rotatingmirrorexperiment}
  Primary-image incidence data, kept separate from the governing laws
\end{definition}
\begin{proof}
  This is the declared model definition of Matches Rotating Mirror Figure.
\end{proof}
\begin{definition}[Has Physical Parameters]
  \label{def:physics:phyx-mini-0025:phyxminiproblems-problemphyxmini0025-hasphysicalparameters}
  \lean{PhyXMiniProblems.ProblemPhyXMini0025.HasPhysicalParameters}
  \uses{def:physics:phyx-mini-0025:phyxminiproblems-problemphyxmini0025-rotatingmirrorexperiment}
  Positivity assumptions selecting the counterclockwise physical branch
\end{definition}
\begin{proof}
  This is the declared model definition of Has Physical Parameters.
\end{proof}
\begin{definition}[Rotating Mirror Optics Laws]
  \label{def:physics:phyx-mini-0025:phyxminiproblems-problemphyxmini0025-rotatingmirroropticslaws}
  \lean{PhyXMiniProblems.ProblemPhyXMini0025.RotatingMirrorOpticsLaws}
  \uses{def:physics:phyx-mini-0025:phyxminiproblems-problemphyxmini0025-dimtime, def:physics:phyx-mini-0025:phyxminiproblems-problemphyxmini0025-pointonbeam, def:physics:phyx-mini-0025:phyxminiproblems-problemphyxmini0025-rotatingmirrorexperiment}
  The physical laws used in the model. the mirror angle advances at constant angular speed; specular reflection makes the outgoing-ray angle advance through twice the mirror's angle because the incident beam is fixed; while the spot is on the east wall, central-ray geometry gives x = (L/2) tan φ; differentiating that relation together with φ' = 2ω gives v = L ω / cos² φ. None of these fields states the requested elapsed-time formula
\end{definition}
\begin{proof}
  This is the declared model definition of Rotating Mirror Optics Laws.
\end{proof}
\begin{definition}[In Closed Time Interval]
  \label{def:physics:phyx-mini-0025:phyxminiproblems-problemphyxmini0025-inclosedtimeinterval}
  \lean{PhyXMiniProblems.ProblemPhyXMini0025.InClosedTimeInterval}
  \uses{def:physics:phyx-mini-0025:phyxminiproblems-problemphyxmini0025-dimtime}
  Membership in a closed physical time interval, expressed in SI seconds
\end{definition}
\begin{proof}
  This is the declared model definition of In Closed Time Interval.
\end{proof}
\begin{definition}[Is East Wall Sweep From OTo Corner]
  \label{def:physics:phyx-mini-0025:phyxminiproblems-problemphyxmini0025-iseastwallsweepfromotocorner}
  \lean{PhyXMiniProblems.ProblemPhyXMini0025.IsEastWallSweepFromOToCorner}
  \uses{def:physics:phyx-mini-0025:phyxminiproblems-problemphyxmini0025-dimtime, def:physics:phyx-mini-0025:phyxminiproblems-problemphyxmini0025-rotatingmirrorexperiment, def:physics:phyx-mini-0025:phyxminiproblems-problemphyxmini0025-inclosedtimeinterval}
  The particular east-wall pass shown in the image: the spot moves from O to the northeast corner without leaving the east wall, and its unwrapped outgoing-ray angle remains on the first-quadrant branch
\end{definition}
\begin{proof}
  This is the declared model definition of Is East Wall Sweep From OTo Corner.
\end{proof}
\begin{definition}[Is Minimum Spot Speed On]
  \label{def:physics:phyx-mini-0025:phyxminiproblems-problemphyxmini0025-isminimumspotspeedon}
  \lean{PhyXMiniProblems.ProblemPhyXMini0025.IsMinimumSpotSpeedOn}
  \uses{def:physics:phyx-mini-0025:phyxminiproblems-problemphyxmini0025-dimtime, def:physics:phyx-mini-0025:phyxminiproblems-problemphyxmini0025-rotatingmirrorexperiment, def:physics:phyx-mini-0025:phyxminiproblems-problemphyxmini0025-inclosedtimeinterval}
  A time realizes the minimum spot speed on the selected wall sweep
\end{definition}
\begin{proof}
  This is the declared model definition of Is Minimum Spot Speed On.
\end{proof}
\begin{definition}[Is Maximum Spot Speed On]
  \label{def:physics:phyx-mini-0025:phyxminiproblems-problemphyxmini0025-ismaximumspotspeedon}
  \lean{PhyXMiniProblems.ProblemPhyXMini0025.IsMaximumSpotSpeedOn}
  \uses{def:physics:phyx-mini-0025:phyxminiproblems-problemphyxmini0025-dimtime, def:physics:phyx-mini-0025:phyxminiproblems-problemphyxmini0025-rotatingmirrorexperiment, def:physics:phyx-mini-0025:phyxminiproblems-problemphyxmini0025-inclosedtimeinterval}
  A time realizes the maximum spot speed on the selected wall sweep
\end{definition}
\begin{proof}
  This is the declared model definition of Is Maximum Spot Speed On.
\end{proof}
% --- Archon declaration topology end ---
% --- Archon physics formalization source end ---
